% Figure: the two condensation modes on a vertical cooled wall, filmwise on the
% left where a continuous liquid film blankets the surface and adds its own
% conduction resistance, and dropwise on the right where discrete droplets grow
% and roll off, leaving bare metal that carries a far higher coefficient.
\begin{figure}[htbp]
\centering
\begin{tikzpicture}[>=stealth]
% filmwise panel
\begin{scope}[shift={(0,0)}]
  % cooled wall
  \fill[enggray!25] (0,0) rectangle (0.35,5);
  \draw[thick] (0.35,0) -- (0.35,5);
  \node[enggray, font=\scriptsize, rotate=90, anchor=south] at (0.15,2.5) {cooled wall};
  % continuous liquid film, widening downward
  \fill[engblue!25] (0.35,0) -- (0.35,5) -- (0.95,5) .. controls (1.25,3) .. (1.35,0) -- cycle;
  \draw[engblue, thick] (0.95,5) .. controls (1.25,3) .. (1.35,0);
  \node[engblue, font=\scriptsize, anchor=west] at (1.4,3.6) {liquid film};
  % downward flow arrows in the film
  \draw[engblue, ->] (0.75,4.2) -- (0.75,3.2);
  \draw[engblue, ->] (0.95,2.5) -- (0.95,1.5);
  \node[font=\scriptsize, anchor=north] at (0.85,-0.1) {filmwise};
  \node[engred, font=\scriptsize, anchor=north] at (0.85,-0.55) {low $h$};
\end{scope}
% dropwise panel
\begin{scope}[shift={(5,0)}]
  \fill[enggray!25] (0,0) rectangle (0.35,5);
  \draw[thick] (0.35,0) -- (0.35,5);
  \node[enggray, font=\scriptsize, rotate=90, anchor=south] at (0.15,2.5) {cooled wall};
  % discrete droplets of varying size
  \foreach \y/\r in {4.6/0.12, 4.0/0.20, 3.3/0.15, 2.7/0.28, 2.0/0.16, 1.4/0.34, 0.7/0.22, 0.3/0.13} {
    \fill[engblue!35] (0.35+\r,\y) circle (\r);
    \draw[engblue, thick] (0.35+\r,\y) circle (\r);
  }
  % a large drop rolling off at the bottom
  \draw[engblue, ->] (1.0,0.9) -- (1.0,0.15);
  \node[engblue, font=\scriptsize, anchor=west] at (1.35,3.8) {droplets};
  \node[font=\scriptsize, anchor=north] at (0.85,-0.1) {dropwise};
  \node[englime!60!black, font=\scriptsize, anchor=north] at (0.85,-0.55) {high $h$};
\end{scope}
\end{tikzpicture}
\caption[Filmwise and dropwise condensation]{The two modes by which a vapour
condenses on a cooled wall. In filmwise condensation a continuous liquid film
covers the surface and thickens as it drains downward, and every unit of heat
must conduct across that film before it reaches the metal, so the film is a
series resistance that caps the coefficient. In dropwise condensation the liquid
beads into discrete droplets that grow, coalesce, and roll off under gravity,
leaving stretches of bare metal that condense vapour directly at a coefficient
several times the filmwise value. Dropwise condensation is the higher-performing
mode, but it depends on a non-wetting surface treatment that erodes in service,
so design coefficients are quoted for the filmwise mode that a real condenser
settles into once the treatment is gone.}
\label{fig:vol04ch06:condensation-modes}
\end{figure}
